% Tabla de abreviaturas y acrónimos generada con el paquete glossaries.
% El estilo de la tabla es long-booktabs

% Definiciones de acrónimos (paquete glossaries).
% - Añadir aquí cualquier sigla nueva; makeindex las ordena alfabéticamente.
% - En el texto puede usarse \gls{clave} (primera vez: forma larga + sigla) o escribir la sigla literal
% - \glsaddall en 6_acronimos.tex garantiza que todas aparezcan en la tabla aunque no se usen con \gls.

\newacronym{api}{API}{\emph{Application Programming Interface}}
\newacronym{bm25}{BM25}{\emph{Best Matching 25}}
\newacronym{etsii}{ETSII}{Escuela Técnica Superior de Ingenieros Informáticos}
\newacronym{etsisi}{ETSISI}{Escuela Técnica Superior de Ingeniería de Sistemas Informáticos}
\newacronym{etsist}{ETSIST}{Escuela Técnica Superior de Ingeniería y Sistemas de Telecomunicación}
\newacronym{etsit}{ETSIT}{Escuela Técnica Superior de Ingenieros de Telecomunicación}
\newacronym{html}{HTML}{\emph{HyperText Markup Language}}
\newacronym{ia}{IA}{Inteligencia Artificial}
\newacronym{idf}{IDF}{\emph{Inverse Document Frequency}}
\newacronym{json}{JSON}{\emph{JavaScript Object Notation}}
\newacronym[longplural={Modelos de lenguaje de gran tamaño}]{llm}{LLM}{Modelo de lenguaje de gran tamaño}
\newacronym{maadm}{MAADM}{Máster Oficial en Aprendizaje Automático y Datos Masivos}
\newacronym{mrr}{MRR}{\emph{Mean Reciprocal Rank}}
\newacronym{pdi}{PDI}{Personal Docente e Investigador}
\newacronym{pie}{PIE}{Proyecto de Innovación Educativa}
\newacronym{rag}{RAG}{\emph{Retrieval-Augmented Generation}}
\newacronym{ri}{RI}{Recuperación de Información}
\newacronym{rrf}{RRF}{\emph{Reciprocal Rank Fusion}}
\newacronym{tfg}{TFG}{Trabajo de Fin de Grado}
\newacronym{tfm}{TFM}{Trabajo de Fin de Máster}
\newacronym{upm}{UPM}{\University}


\glsaddall  % incluye todos los acrónimos definidos, se usen o no con \gls

\printglossary[type=\acronymtype, style=long-booktabs,
               title={Abreviaturas y acrónimos},
               toctitle={Abreviaturas y acrónimos}]
